\section{Mixed connection values}\label{sec:gamma}
All nodes in this section are positive algebraic numbers; all logarithms
are real. Put $\mathcal E=\operatorname{acl}_{\C}\Qbar(e^\beta:\beta\in\Qbar)$.
Write $R_B(u)=\mathsf S(u)/I_0(u)$ and $R_a(t)=S_a(t)/f_a(t)$ for the
argument-normalized functions of Section~\ref{sec:setup}.
For $a\in\Q\setminus\Z$, write $a=\vartheta_a+k_a$, where
$\vartheta_a=p_a/q_a\in(0,1)$ is reduced and $k_a\in\Z$. Define
\[
 G_\vartheta=\psi(\vartheta)+\gamma,\qquad
 H(\vartheta,k)=\begin{cases}
 \sum_{j=0}^{k-1}(\vartheta+j)^{-1},&k\ge0,\\
 -\sum_{j=k}^{-1}(\vartheta+j)^{-1},&k<0.
 \end{cases}
\]
\begin{theorem}[Mixed connection values]\label{thm:connection}\label{ct2:statement-start}
Select $N\ge1$ rows of the following table, with positive algebraic
nodes $x,u,t$, integer orders $m\ge0$ and rational nonintegral parameters $a$.
Require distinct Ein nodes and distinct Bessel nodes.
Exclude Kummer pairs at the same $t$ whose parameters differ by an integer.
Also exclude Bessel--Kummer pairs with $t=2u$ and $a\in\tfrac12+\Z$.
Put $\epsilon_m=(-1)^m$ and define
\[
\begin{gathered}
 (p_0,v_0)=(1,0),\quad(p_1,v_1)=(0,1),\\
 (p_{m+1},v_{m+1})=(p_{m-1},v_{m-1})-\frac{2m}{u}(p_m,v_m),
 \quad\beta_m=-\frac{\epsilon_mv_m}{uI_0(u)I_m(u)}.
\end{gathered}
\]
The denominators below are nonzero. The table defines $V=X-d$:
\begin{equation}\label{eq:ct2-table}
\begin{array}{c|c|c}
 V_i&X_i&d_i\\ \hline
 E_1(x)&\Ein(x)&\gamma+\log x\\
 K_m(u)/I_m(u)&\epsilon_mR_B(u)+\beta_m(u)&
       \epsilon_m\{\gamma+\log(u/2)\}\\
 \Gamma(a)U(a,1,t)/f_a(t)&R_a(t)-H(\vartheta_a,k_a)&
       \gamma+\log t+G_{\vartheta_a}
\end{array}
\end{equation}
For an arbitrary integer Bessel order, use its absolute value. Set
\[
 B=\operatorname{acl}_{\C}\mathcal E
 \bigl(I_0(u),I_1(u),f_a(t),f'_a(t):\text{selected rows}\bigr).
\]
Gauss's formula and multiplicative relations among the positive algebraic
arguments of the logarithms express $d=A\mathbf c$, with $A\in M_{N,r_{\mathbf c}}(\Qbar)$
and $r_{\mathbf c}$ the number of retained constants.
Here $\mathbf c$ consists of $\gamma$, the required real logarithms, and $\pi$
only if its coefficient is nonzero; logarithms are combined and unused entries omitted.
For each $F=\mathcal E,B$,
\begin{equation}\label{eq:ct2-actual-degree}
 \trdeg_F F(\mathbf c,V_1,\ldots,V_N)\ge N.
\end{equation}
Thus these numbers contain an algebraically independent subset of
size $N$ over either base field. If all entries of $\mathbf c$ are algebraic over
$F$, the values $V_1,\ldots,V_N$ themselves are algebraically independent
 over $F$.
\label{ct2:statement-end}
\end{theorem}
\begin{proof}\label{ct2:proof-start}
The connection expansions \cite[6.2.4, 10.28.2, 10.31.1, 13.2.9]{DLMF} give
\begin{align*}
 E_1(x)&=\Ein(x)-\gamma-\log x,\\
 K_m(u)/I_m(u)&=\epsilon_m\{R_B(u)-\gamma-\log(u/2)\}+\beta_m(u),\\
 \Gamma(a)U(a,1,t)&=S_a(t)-f_a(t)\{\log t+\psi(a)+2\gamma\}.
\end{align*}
The Bessel identity follows from $I_m=p_mI_0+v_mI_1$,
$\epsilon_mK_m=p_mK_0-v_mK_1$ and $I_0K_1+I_1K_0=1/u$.
In particular $I_m(u)>0$ and $\beta_m(u)\in B$.
Theorem~\ref{thm:values}, applied to the closed normalized Kummer block
at $z=1$, gives $f_a(t)\ne0$, including for negative rational $a$.
This single-block class is its own representative, so its regular
coordinate is $f_a(t)$ in \eqref{eq:value-nonzero}.
The recurrence gives $\psi(a)+\gamma=G_{\vartheta_a}+H(\vartheta_a,k_a)$.
Gauss's formula is
\begin{equation}\label{eq:ct2-gauss}
 G_{p/q}=-\log(2q)-\frac\pi2\cot\frac{\pi p}{q}
 +2\sum_{j=1}^{\lfloor(q-1)/2\rfloor}
       \cos\frac{2\pi pj}{q}\log\sin\frac{\pi j}{q}.
\end{equation}
Its coefficients and positive logarithm arguments are algebraic.
Taking a multiplicative $\Q$-basis of the arguments proves $d=A\mathbf c$
with the stated reduced constant column.

The exclusions place the selected second-order blocks in distinct linkage classes.
Distinct Ein nodes have independent residues $e_{(0,0)}-e_{(-x,0)}$.
Apply closure, normalization and rate saturation.
By Theorem~\ref{thm:values} and Corollary~\ref{cor:value-base}, the Ein values and one triple
$(I_0,I_1,\mathsf S)$ or $(f_a,f'_a,S_a)$ per selected class are then jointly independent
 over $\mathcal E$. Adjoin the regular pairs and divide by their nonzero
regular coordinates. Then pass to the algebraic closure $B$.
The remaining $N$ ratios are independent. Nonzero diagonal scalings and
$B$-valued shifts show that $X_1,\ldots,X_N$ are independent over $B$,
hence over $\mathcal E$.
Since $X=V+A\mathbf c$ and $A$ is algebraic, $F(X)\subset F(\mathbf c,V)$ proves
\eqref{eq:ct2-actual-degree}. A finite generating set contains a
transcendence basis. If $\mathbf c$ is algebraic over $F$, it lies in $F$,
since both permitted base fields are algebraically closed within $\C$.
Then $V=X-A\mathbf c$ is an $F$-valued translation of $X$.
\label{ct2:proof-end}
\end{proof}

\begin{corollary}\label{cor:gamma-gompertz}
At least one of $\gamma$ and $E_1(1)$, equivalently at least one of
$\gamma$ and $\delta=eE_1(1)=\int_0^\infty e^{-s}/(1+s)\,ds$,
is transcendental over $\mathcal E$.
\end{corollary}
\begin{proof}
Use the Ein row at $x=1$: $B=\mathcal E$, $\mathbf c=(\gamma)$ and
$\Ein(1)=\gamma+E_1(1)$. Also $e\in\mathcal E^\times$.
\end{proof}
The stated alternative over the full field $\mathcal E$ is known.
It follows from Shidlovskii's Chapter~7, \S1, Theorem~1
\cite[pp.~232--233]{Shidlovskii1989}, applied after the finite-rate
reduction in Supplementary Section~\ref{sup:classical-comparison}.
Mahler announced an exponential-integral statement without proof
\cite[\S23]{Mahler1968}. Rivoal proves the one-point result for
algebraic $x\notin(-\infty,0]$ \cite[Theorem~2(ii), \S7]{Rivoal2012}.
Theorem~\ref{thm:connection} assembles multi-node mixed families over $B$.
The comparison with the sources is partial.

\begin{example}\label{ex:ct2-mixed}
Choose $x=1$, $(u,m)=(2,0)$ and $(t,a)=(1,1/2)$.
The rows are admissible since the mixed linkage would require $t=4$.
Here $\mathbf c=(\gamma,\log2)^t$ and
\[
 B_*=\operatorname{acl}_{\C}\mathcal E
 (I_0(2),I_1(2),f_{1/2}(1),f'_{1/2}(1)).
\]
For each $F=\mathcal E,B_*$, the five numbers
\[
 \gamma,\quad\log2,\quad E_1(1),\quad
 \frac{K_0(2)}{I_0(2)},\quad
 \frac{\Gamma(1/2)U(1/2,1,1)}{f_{1/2}(1)}
\]
contain three algebraically independent numbers over $F$.
Supplementary Section~\ref{sup:cancellation-remarks} computes $A$ and an explicit cancellation.
\end{example}

\begin{proposition}[Coefficient cancellation]\label{prop:coefficient-cancellation}\label{lem:coefficient-cancellation}
Let $\kappa$ be the dimension of $\Span_B\{d_i+B\}$ in $\C/B$.
Let $J,J_0$ have row bases of $\{\lambda:\lambda d\in B\}$ and
$\{\lambda:\lambda d=0\}$ in $B^{1\times N}$.
For independent indeterminates $Z$, put
$\mathcal A_d=\{P\in B[Z]:P(Z-d)\in B[Z]\}$. Then
\begin{equation}\label{eq:ct2-algebra}
\begin{gathered}
 \mathcal A_d=B[JZ],\qquad
 \{P\in B[Z]:P(Z-d)=P(Z)\}=B[J_0Z],\\
 B[JV]=B[JX],\qquad\trdeg_B B(JV)=N-\kappa,\qquad
 N-\rank_{\Qbar}A\le N-\kappa\le N.
\end{gathered}
\end{equation}
\end{proposition}
Supplementary Section~\ref{sup:cancellation-remarks} proves the proposition in full.
Neither $\kappa$ nor $J$ is claimed effective. Even whether $\gamma+\log x\in B$ is undecided.
Accidental intersections $B[V]\cap B[X]$ and equalities $P(V)=P(X)$ are not classified.
No independence of $\mathbf c$ or of $X$ over $B(\mathbf c)$ is assumed.
Linked rows must be reduced first; higher Ein tails are excluded.
The irrationality of Euler's constant remains open.
\label{ct2:section-end}
